\section{Example 1.1}\label{sec:example1-1}

\subsection{Problem description}\label{subsec:example1-1-problem}

A decision maker has three alternatives, A, B and C, and three possible states of nature \(S_1\), \(S_2\) and \(S_3\).
State probabilities are \(P(S_1)=0.30\), \(P(S_2)=0.40\) and \(P(S_3)=0.30\).
The payoffs (in generic units) are:

\begin{itemize}
	\item Alternative A: \(S_1=150\), \(S_2=80\), \(S_3=-20\).
	\item Alternative B: \(S_1=110\), \(S_2=90\), \(S_3=40\).
	\item Alternative C: \(S_1=70\), \(S_2=70\), \(S_3=60\).
\end{itemize}

\subsection{Solution}\label{subsec:example1-1-solution}

\subsubsection{(C2) Interpreting the parts of the decision problem}\label{subsubsec:example1-1-c2}

We identify the components of the decision situation.

\begin{center}
	\begin{tabularx}{\linewidth}{@{}L{0.35\linewidth}Y@{}}
		\toprule
		Element & Description \\
		\midrule
		Decision alternatives & Alternatives A, B, C \\
		States of nature & \(S_1\), \(S_2\), \(S_3\) \\
		Fortuitous events & The state of nature that occurs, unknown at the time of the decision \\
		Consequences & Payoff (generic units) for each alternative--state pair \\
		Probabilities & \(P(S_1)=0.30\), \(P(S_2)=0.40\), \(P(S_3)=0.30\) \\
		\bottomrule
	\end{tabularx}
\end{center}

\subsubsection{(C3) Building the payoff table from the description}\label{subsubsec:example1-1-c3}

From the description we build the payoff table.

\begin{center}
	\begin{tabularx}{\linewidth}{@{}L{0.25\linewidth}C C C C@{}}
		\toprule
		State of nature & Probability & Alternative A & Alternative B & Alternative C \\
		\midrule
		\(S_1\) & 0.30 & 150 & 110 & 70 \\
		\(S_2\) & 0.40 & 80 & 90 & 70 \\
		\(S_3\) & 0.30 & -20 & 40 & 60 \\
		\bottomrule
	\end{tabularx}
\end{center}

\subsubsection{(C4) Maximax rule decision making without probabilities}\label{subsubsec:example1-1-c4}

We ignore probabilities and focus on the best payoff for each alternative.

\[
\begin{aligned}
\max_i \text{Payoff}(A, S_i) &= 150\\
\max_i \text{Payoff}(B, S_i) &= 110\\
\max_i \text{Payoff}(C, S_i) &= 70
\end{aligned}
\]
\[
\max\{150, 110, 70\} = 150
\]
The largest of these is 150, so the Maximax rule recommends Alternative A.

\subsubsection{(C5) Maximin rule decision making without probabilities}\label{subsubsec:example1-1-c5}

We focus on the worst payoff for each alternative.

\[
\begin{aligned}
\min_i \text{Payoff}(A, S_i) &= -20\\
\min_i \text{Payoff}(B, S_i) &= 40\\
\min_i \text{Payoff}(C, S_i) &= 60
\end{aligned}
\]
We then select the highest of these worst outcomes:
\[
\max\{-20, 40, 60\} = 60
\]
The Maximin rule recommends Alternative C.

\subsubsection{(C6) Maximum opportunity criterion minimax regret}\label{subsubsec:example1-1-c6}

We compute the regret table by comparing each payoff to the best payoff in each state.

Best payoff in each state
\[
\begin{aligned}
\text{S1} &:\ \max\{150, 110, 70\} = 150\\
\text{S2} &:\ \max\{80, 90, 70\} = 90\\
\text{S3} &:\ \max\{-20, 40, 60\} = 60
\end{aligned}
\]

Regrets are best payoff minus actual payoff.

\begin{center}
	\begin{tabularx}{\linewidth}{@{}L{0.25\linewidth}C C C@{}}
		\toprule
		State of Nature & Alternative A & Alternative B & Alternative C \\
		\midrule
		\(S_1\) & \(150-150=0\) & \(150-110=40\) & \(150-70=80\) \\
		\(S_2\) & \(90-80=10\) & \(90-90=0\) & \(90-70=20\) \\
		\(S_3\) & \(60-(-20)=80\) & \(60-40=20\) & \(60-60=0\) \\
		\bottomrule
	\end{tabularx}
\end{center}

Maximum regret for each alternative
\[
\begin{aligned}
\max_i \text{Regret}(A, S_i) &= \max\{0, 10, 80\} = 80\\
\max_i \text{Regret}(B, S_i) &= \max\{40, 0, 20\} = 40\\
\max_i \text{Regret}(C, S_i) &= \max\{80, 20, 0\} = 80
\end{aligned}
\]
\[
\min\{80, 40, 80\} = 40
\]
The minimax regret rule chooses Alternative B.

\subsubsection{(C7) Using probabilities and expected value Bayes decision rule}\label{subsubsec:example1-1-c7}

We compute the expected monetary value (EMV) for each alternative.

For Alternative A
\[
\begin{aligned}
EMV_A &= 0.30(150) + 0.40(80) + 0.30(-20) \\
&= 45 + 32 - 6 = 71
\end{aligned}
\]

For Alternative B
\[
\begin{aligned}
EMV_B &= 0.30(110) + 0.40(90) + 0.30(40) \\
&= 33 + 36 + 12 = 81
\end{aligned}
\]

For Alternative C
\[
\begin{aligned}
EMV_C &= 0.30(70) + 0.40(70) + 0.30(60) \\
&= 21 + 28 + 18 = 67
\end{aligned}
\]
\[
\max\{71, 81, 67\} = 81
\]
Under the Bayes EMV rule, the highest EMV is 81, so the rule recommends Alternative B.

\subsubsection{(C8) Comparing Bayes, Maximin, Maximax and minimax regret}\label{subsubsec:example1-1-c8}

Summary of recommendations from the different rules

\begin{itemize}
	\item Bayes EMV recommends Alternative B.
	\item Minimax regret recommends Alternative B.
	\item Maximax recommends Alternative A.
	\item Maximin recommends Alternative C.
\end{itemize}

The EMV and minimax regret rules both favor the same alternative, while Maximax and Maximin highlight different extremes of the payoff table.

\subsubsection{(C9) Tree diagram of the decision-making problem}\label{subsubsec:example1-1-c9}

We represent the same situation as a decision tree.

This diagram summarises the possible states and payoffs for each alternative.

\subsubsection{(C10) Evaluating all possibilities in the decision problem}\label{subsubsec:example1-1-c10}

\begin{itemize}
	\item Alternative A maximises the best payoff but performs poorly in worst-case and regret terms.
	\item Alternative C protects against the worst case but yields lower expected value.
	\item Alternative B is favoured by both expected value and minimax regret.
\end{itemize}

Overall, the comparison shows how the choice depends on risk attitude, with Alternative B providing a balanced option.

\subsubsection{(C1) Formulating the final decision and summarising all criteria}\label{subsubsec:example1-1-c1}

Summary of final recommendations across all decision criteria

\begin{itemize}
	\item Maximax favors Alternative A, prioritising the highest possible payoff.
	\item Maximin favors Alternative C, emphasising protection against the worst case.
	\item Minimax regret favors Alternative B, limiting potential regret.
	\item Bayes EMV favors Alternative B, maximising expected value.
\end{itemize}

Considering all criteria together, Alternative B represents the most coherent final decision. It is supported by both expected value and minimax regret and avoids extreme downside outcomes.

\vspace{6pt}
\noindent\hyperref[table-of-contents]{Back to Top}
